\section{从发布会到真实效果：模型公司声量的变化}

张鹏在后半段提到一个很微妙的变化：过去模型发布常常需要隆重发布会、几百人现场、强宣发；但 GLM-4.5 开源后，智谱并没有做特别大的发布，海外开发者和产品却给了很高评价，甚至出现 Windsurf 等产品替换或接入模型、其他厂商蒸馏和裁剪模型的现象。这个变化说明，大模型公司的技术声量正在从“发布会叙事”转向“真实使用效果”。

\subsection{为什么发布会边际收益下降}

前面说模型声量转向真实效果，本节解释原因。模型发布太多以后，叙事本身不再稀缺，真实使用和成本效率才更稀缺。

大模型行业经历多轮发布后，用户和开发者对“又一个模型发布”已经审美疲劳。宣传可以带来短期注意力，但模型能否被开发者接入、被客户买单、被社区自发讨论，越来越取决于真实能力和成本效率。DeepSeek 的影响，也在于它把“少说、多开源、效果说话”的路径推到行业前台。

\begin{knowledgebox}{模型声量的三种来源}
\begin{tabularx}{\textwidth}{p{0.22\textwidth}p{0.34\textwidth}X}
\toprule
来源 & 优点 & 风险 \\
\midrule
发布会 & 叙事集中、容易制造里程碑 & 用户审美疲劳，难证明长期价值。 \\
Benchmark & 可比较、可传播、便于技术判断 & 容易过拟合榜单，和真实任务有距离。 \\
真实使用 & 开发者、客户和产品自然反馈 & 反馈慢、不可控，但最接近商业价值。 \\
\bottomrule
\end{tabularx}
\end{knowledgebox}

\subsection{被使用也是一种评测}

当一个模型被开发者接入，被产品替换原有模型，被其他团队蒸馏或二次包装，它已经不只是 paper 或榜单里的模型，而是进入了实际生态。这里有好处，也有麻烦：好处是模型能力被真实任务验证；麻烦是价值捕获变复杂，开源影响不一定直接变成收入。

\begin{importantbox}{真实效果的评测逻辑}
模型公司不能只问“我们榜单排第几”，还要问“谁在用，为什么用，替代了谁，成本是否更低，客户是否愿意持续付费”。这才是上市公司阶段更重要的能力证明。
\end{importantbox}

\subsection{本章小结}

模型发布的中心正在从“我说我很强”转向“别人用我做成了什么”。这也是开源生态、商业化和模型能力之间的新关系。

